\documentclass[10pt,a4paper]{amsart}
\usepackage[margin=19mm]{geometry}
\usepackage{amsmath,amssymb,mathtools}
\usepackage[scr=rsfs,cal=euler]{mathalfa}
\usepackage{xfrac}
\usepackage[unicode,hidelinks,bookmarks=false]{hyperref}
\newcommand{\ie}{i.e.\ }
\newcommand{\cf}{cf.\ }
\newcommand{\resp}{respectively\ }
\newcommand{\Subsection}{\subsection}
\providecommand{\nobreakdash}{\nobreak\hbox{-}}
\setlength{\emergencystretch}{2em}
\multlinegap0pt
\usepackage{amssymb}
\usepackage{xcolor}
\usepackage{enumerate}
\usepackage{tikz}
\usepackage[shortlabels]{enumitem}
\usepackage{mathtools}
\usepackage{bm}
\newtheorem{lemma}{Lemma}
\newtheorem{theorem}{Theorem}
\newcommand{\CK}[1]{\color{red}{#1}\color{black}}
\def \F {\mathbf{F}}
\def \M {\mathcal{M}}
\def \R {\mathbb{R}}
\def \c{c_{k_1,k_2,r,d}}
\def \dis {\displaystyle}
\title{Monotonicity of the propagation speed with respect to the diffusion in a Lotka-Volterra competition-diffusion system}
\author{Cyrille Kenne}
\date{Source: arXiv:2608.27431v1 (CC BY 4.0)}
\begin{document}
\maketitle

\noindent\emph{Source and licence.} Monotonicity of the propagation speed with respect to the diffusion in a Lotka-Volterra competition-diffusion system. Version arXiv:2608.27431v1 (CC BY 4.0), distributed by arXiv under the Creative Commons Attribution 4.0 International licence, http://creativecommons.org/licenses/by/4.0/, as recorded in the arXiv licence metadata for this exact version and shown on the abstract page https://arxiv.org/abs/2608.27431v1. Full licence notice: \texttt{licenses/arxiv-2608.27431-CC-BY-4.0.txt}.\par\smallskip

\noindent\textit{Editorial scope.} Counted unit: the article's theorem theoimportant (source0.tex lines 744--753). The article's complete original proof of it is delivered with it (source0.tex lines 754--814, one \texttt{proof} environment of 61 lines). No mathematics is written, repaired or re-derived: the statement and the proof are the article's own lines, in the article's own notation, and no retained line is substituted or re-flowed.  Retained uncounted as the source-local context the proof needs: lines 318--334; lines 718--725; lines 727--734; lines 736--738.  Nothing was omitted from any retained span, so the package carries no omission record.  No retained line contains a citation, so the wrapper carries no bibliography and no bibliography entry.  No retained line contains a figure environment, an image-inclusion command or drawing code, so no image is delivered and the package declares no asset dependency.  Editorial formatting only, in addition: an AMS \texttt{amsart} preamble that restates the article's own declarations and macros used by the retained lines, namely the environments \texttt{lemma}, \texttt{theorem} on separate counters, together with the standard packages those lines need.  The retained environments print in this document as Lemma 1, Theorem 1 in order of appearance, whereas the article prints the same items with the numbering of its own full document.  The complete native source of the article as distributed in the arXiv e-print for this exact version is bundled unchanged as \texttt{sources/208714/source0.tex}.
\begin{lemma}\label{lemdecay} Let $\pi_0=(k_{1,0}, k_{2,0},r_0,d_0) \in Q$ and let $(\Phi_0, c_0)$ be its monotone front. There exist $C,\eta>0$ such that
	\begin{equation}\label{cy0}
		\|\Phi_0(\xi)-E_-\| + \|\Phi_0'(\xi)\|  \leq C e^{\eta\xi}, \qquad \xi\leq 0
	\end{equation}
	and 
	\begin{equation}\label{cy1}
		\|\Phi_0(\xi)-E_+\| + \|\Phi_0'(\xi)\|  \leq C e^{-\eta\xi}, \qquad \xi\geq 0.
	\end{equation}
	In particular,
	\begin{equation}\label{cy2}
		\Phi'_0\in H^2(\R)^2, \qquad \Phi_0'' \in L^2(\R)^2, \qquad \F_{k_{1,0}, k_{2,0},r_0}(\Phi_0)\in L^2(\R)^2.
	\end{equation}
Moreover, for every integer $m\geq 1$, there exists $C_m>0$ such that
	\begin{equation}\label{generaldecay}
		\|\Phi_0^{(m)}(\xi)\|\leq C_me^{-\eta|\xi|}, \qquad \text{ for all } \xi\in \R.
	\end{equation}
\end{lemma}

\begin{equation}\label{syst2new}
	\left\{ 
	\begin{array}{llllll}
		p''+cp'+(1-2U-k_1V)p+k_1Uq=0,\\
		dq''+cq'+r(1-2V-k_2U)q+rk_2Vp=0.
	\end{array} 
	\right.
\end{equation}

\begin{equation}\label{systab}
	\left\{ 
	\begin{array}{llllll}
		a''-ca'+(1-2U-k_1V)a+rk_2Vb=0,\\
		db''-cb'+r(1-2V-k_2U)b+k_1Ua=0,
	\end{array} 
	\right.
\end{equation}

\begin{equation}\label{limab}
(a,b)(\pm \infty)=0.
\end{equation}

\begin{theorem}\label{theoimportant}
	Let $k_1,k_2>1$, $r>0$ and  $c:=\c$. We assume that 
	\begin{equation}\label{condition}
	c\leq 0	 \; \text{ and }\; c\left(1-\frac{1}{d}\right)\leq 0
	\end{equation}
hold. Then
	\begin{equation}\label{eq04}
		K:=\int_{\R}q'(\xi)b(\xi)d\xi>0.
	\end{equation}
\end{theorem}

\begin{proof}
We first observe that an integration by parts over $\R$ while using Lemma \ref{lemdecay} and \eqref{limab} gives
\begin{equation}\label{important0}
	2K= \int_{\R}\left(q'(\xi)b(\xi)-q(\xi)b'(\xi) \right)d\xi. 
\end{equation}
Now, we multiply the first equation in \eqref{syst2new} by $a$ and the first equation in \eqref{systab} by $p$ and we subtract the resulting equations to obtain
\begin{equation*}
a''p-p''a-c(p'a+a'p)=k_1Uqa-rk_2Vbp.
\end{equation*}
since $(a'p-p'a)'=a''p-p''a$, we can rewrite the previous equation as 
\begin{equation}\label{siha0}
	(a'p-p'a-cap)'=k_1Uqa-rk_2Vbp.
\end{equation}
Next, multiplying the second equation in \eqref{syst2new} by $b$ and the second equation in \eqref{systab} by $q$ and combining the resulting equations, we obtain
\begin{equation}\label{siha1}
	[d(q'b-b'q)+cbq]'=k_1Uqa-rk_2Vbp.
\end{equation}
On the one hand, combining \eqref{siha0} and \eqref{siha1} leads us to 
\begin{equation*}
	[d(q'b-b'q)+cbq]'=(a'p-p'a-cap)'.
\end{equation*}
So, the functions $\xi \mapsto d(q'(\xi)b(\xi)-b'(\xi)q(\xi))+cb(\xi)q(\xi)$ and $\xi \mapsto a'(\xi)p(\xi)-p'(\xi)a(\xi)-ca(\xi)p(\xi)$ have the same derivatives. Using Lemma \ref{lemdecay} and thanks to the fact that $\xi \mapsto a'(\xi)$ and $\xi\mapsto b'(\xi)$ are bounded (because $H^1(\R)\hookrightarrow  C_b(\R)$), we can deduce that these two functions have the same limit $0$ at infinity. Therefore, these two functions must be equal. We set for all $\xi \in \R$,
\begin{equation}\label{defG}
G(\xi):=d(q'(\xi)b(\xi)-b'(\xi)q(\xi))+cb(\xi)q(\xi)=a'(\xi)p(\xi)-p'(\xi)a(\xi)-ca(\xi)p(\xi).
\end{equation} 
Then we have 
\begin{equation}\label{limderivG}
\lim_{\xi\to \pm\infty}G(\xi)=0  \;  \text{ and } \; G'(\xi)= k_1U(\xi)q(\xi)a(\xi)-rk_2V(\xi)b(\xi)p(\xi).
\end{equation}
Differentiating the equation in \eqref{limderivG}, gives
\begin{equation}\label{siha2}
G''=k_1\left[-pqa+ Uq'a+Uqa'\right] -rk_2\left[qbp+Vb'p+Vbp'\right],
\end{equation}
where we dropped the variable $\xi$ for clarity.
From \eqref{defG}, we can write $a'=\frac{G}{p}+\frac{p'}{p}a+ca$ and $q'=\frac{G}{d b}+\frac{b'}{b}q-\frac{c}{d}q$. Substituting these values in \eqref{siha2} lead us to
\begin{equation*}
G''=\left[ \left( \frac{b'}{b}+\frac{p'}{p}\right)G' +k_1U\left( \frac{a}{d b}+\frac{q}{p}\right) G+ k_1c\left(1-\frac{1}{d}\right)Uqa -pq(k_1a+rk_2b)\right].
\end{equation*}
That is 
\begin{equation}\label{siha3}
\M G:=G''-\left( \frac{b'}{b}+\frac{p'}{p}\right)G' -k_1U\left( \frac{a}{d b}+\frac{q}{p}\right) G=q\left[ k_1c\left(1-\frac{1}{d}\right)Ua -p(k_1a+rk_2b)\right].
\end{equation}
Since $k_1,k_2>1, q,U,a,p,r,b>0$,  if $c\left(1-\frac{1}{d}\right)\leq0$, then the right hand side of \eqref{siha3} is strictly negative. %\CK{This condition seems strong!!! Is there a way to weaken it???????????}
That is  
\begin{equation}\label{siha4}
\M G(\xi)<0\; \text{ for all }\;  \xi\in \R.
\end{equation}
Note that the same condition  and \eqref{siha3} imply that $G$ is not identically $0$. Now, we claim that $G>0$ on $\R$. Indeed, suppose that there is $\xi_0\in \R$ so that $G(\xi_0)\leq 0$.  Then because  $\dis \lim_{\xi\to \pm\infty}G(\xi)=0$ and $G$ is continuous, $G$ attains a nonpositive minimum at some point $\xi_1 \in \R$. Thus at $\xi_1$, we have 
\begin{equation*}
G(\xi_1)\leq 0, \quad G'(\xi_1)=0 \text{ and } G''(\xi_1)\geq 0.
\end{equation*}
Since $k_1U\left( \frac{a}{d b}+\frac{q}{p}\right)>0$, we obtain from \eqref{siha3} that $\M G(\xi_1)\geq 0$. This contradicts \eqref{siha4}. Hence $G>0$ on $\R$. Now, from \eqref{defG}, we have 
\begin{equation*}
q'(\xi)b(\xi)-b'(\xi)q(\xi)=\frac{1}{d}\left( G(\xi)-cb(\xi)q(\xi)\right),
\end{equation*}
which combined with \eqref{important0} imply 
\begin{equation*}
K= \frac{1}{2d}\int_{\R}\left( G(\xi)-cb(\xi)q(\xi)\right)d\xi. 
\end{equation*}
Combining $G>0$ and $c\leq 0$, we deduce that $K>0$.  This completes the proof.
\end{proof}

\end{document}
